\scanpage{17}
\mspage{8}
en termes de l'\uline{action extérieure} correspondante de $\Gamma$ sur $\pi$, comme l'image inverse de l'extension universelle
\[
(15)\qquad
\diagram{\begin{tikzcd}[ampersand replacement=\&]
1\arrow[r]\&\pi\arrow[r]\&\operatorname{Aut}(\pi)\arrow[r]\&\operatorname{Aut\,ext}(\pi)\arrow[r]\&1\\
\&\&\&\Gamma\arrow[u,"\varphi"']\&
\end{tikzcd}}
\]

Pour $K$ fixé, donc $k$ fixé, dire qu'on fixe un $\Gamma=\Gamma_{\overline{\mathbb Q}/k}$ revient à dire qu'on fixe une clôture alg. de $k$, dire qu'on fixe \inserted{un} $\pi_{\bar K/K}=\pi_1(K\otimes_k\bar k)$ signifie \inserted{\unsure{est}} qu'on fixe un rev. universel de $\operatorname{Spec}(K\otimes_k\bar k)=\eta\otimes_k\bar k$, \deleted{\unsure{compatible}} les deux ensembles reviennent à se donner le rev. universel \deleted{$\unsure{S\bar K|\bar K}$} $\bar\eta=\operatorname{Spec}(\bar K)$ de $K$. Par la suite, nous décrirons (avec une fidélité qui \unsure{met|me tient} à cœur) les couples $(K,\bar K)$ d'une ext.\inserted{$K$} de $\mathbb Q$ de type fini, et d'une clôture algébrique $\bar K$ de $K$, par les triples $(\pi,\Gamma,\varphi)$, où $\pi=\pi_{\bar K/K}$ \inserted{et $\Gamma=\Gamma_{\bar K/K}$} sont des groupes profinis, et $\varphi:\Gamma\unsure{-|\longrightarrow}\operatorname{Aut\,ext}(\pi)$ une action extérieure de $\Gamma$ sur $\pi$ --- ce qui
\margin{16}
